%%% ------------------------------------------------------------------
%%% KOL/MORS Morpho-Syntax
%%% Assignment: describe a morpho-syntactic phenomenon
%%%
%%% This is a skeleton.  Replace everything marked TODO, delete the
%%% parts you do not need, and keep the parts you do.
%%%
%%% The packages below are the ones recommended in
%%%   "A very short guide to formatting HPSG with LaTeX"
%%%     https://www.overleaf.com/read/hnjksbxgbhtq
%%% which has worked examples of trees, feature structures and glosses.
%%% Read that first if you have not used LaTeX for linguistics before.
%%%
%%% To compile on your own machine:
%%%     pdflatex assignment
%%%     bibtex   assignment
%%%     pdflatex assignment
%%%     pdflatex assignment
%%% (twice at the end, so the citations and cross-references settle)
%%%
%%% On Overleaf just press Recompile; it does all of that for you.
%%% ------------------------------------------------------------------

\documentclass[a4paper,12pt]{article}

%%% uniform margins
\usepackage[margin=25mm]{geometry}

\usepackage[utf8]{inputenc}       % delete this line if you compile with XeLaTeX
\usepackage[T1]{fontenc}          % proper accented letters: ě š č ř ž ý á í é
\usepackage[english]{babel}

%%% citations: \citep{...} and \citet{...}
\usepackage[round]{natbib}
\bibliographystyle{plainnat}

%%% hyperlinks
\usepackage[svgnames]{xcolor}
\usepackage{hyperref,url}
\hypersetup{
    colorlinks=true,
    linkcolor=DarkBlue,
    urlcolor=DarkRed,
}

%%% trees
\usepackage{forest}
\useforestlibrary{linguistics}
\forestapplylibrarydefaults{linguistics}

%%% AVMs (attribute value matrices), from Language Science Press
\usepackage{langsci-avm}

%%% numbered examples and glosses
\usepackage{langsci-gb4e}
\let\eachwordone=\itshape         % italicises the source line of a gloss

%%% Some versions of gb4e need \noautomath so that _ and ^ behave
%%% outside maths; older ones do not have it at all, so provide a
%%% harmless fallback rather than let the file fail to compile.
\providecommand{\noautomath}{}
\noautomath

%%% A shorthand for glossing labels, so they come out as small capitals:
%%% type \gm{acc} and you get ACC.
\newcommand{\gm}[1]{\textsc{\MakeLowercase{#1}}}

\title{TODO: a title that names the phenomenon\\
  \large KOL/MORS Morpho-Syntax}
\author{TODO: First Author's Name \and TODO: Second Author's Name}
\date{2027-01-31}

\begin{document}
\maketitle

\begin{abstract}
  TODO: three or four sentences.  What is the phenomenon, what is
  puzzling about it, and what do you conclude?  Write this last.
\end{abstract}

%%% ------------------------------------------------------------------
\section{Introduction}
%%% ------------------------------------------------------------------

TODO: Say what the phenomenon is and why it is interesting.
``Interesting'' is subjective, so argue for it: what would we expect,
and what happens instead?  Lead with a concrete example rather than
with the literature.

Here is how you give a numbered example:

\begin{exe}
  \ex \label{ex:good} The more you read, the better you write.
\end{exe}

Refer back to it as (\ref{ex:good}) --- never as ``the example above''.

A star marks an unacceptable example.  Grouped examples share one
number and get letters:

\begin{exe}
  \ex \begin{xlist}
    \ex Who did you say left?
    \ex[*]{ Who did you say that left? }
  \end{xlist}
\end{exe}

%%% ------------------------------------------------------------------
\section{The data}
%%% ------------------------------------------------------------------

TODO: Present the examples your argument rests on.  If they are not in
English they must be glossed, so that a reader who does not speak the
language can follow.

This is a glossed example.  The source line and the gloss line must
have the \emph{same number of words}, and matching hyphens, or the
columns will not line up:

\begin{exe}
  \ex Czech [ces] \\
      \gll Pět   žen     přišlo            \\
           five  women   come.\gm{pst}.\gm{n}.\gm{sg} \\
  \trans ``Five women came.''
\end{exe}

\begin{exe}
  \ex Czech [ces] \\
      \gll Nikdo   nic      ne-řekl          \\
           nobody  nothing  \gm{neg}-say.\gm{pst}.\gm{m}.\gm{sg} \\
  \trans ``Nobody said anything.''
\end{exe}

Three conventions, all of which we will mark:

\begin{itemize}
\item Cited forms go in \emph{italics}.  In a gloss the source line is
  italicised for you, but in running text you have to do it yourself:
  the word \emph{žena}, the suffix \emph{-ů}.
\item Glosses of grammatical material go in \gm{small capitals};
  lexical material stays lower case.  Follow the
  \href{https://www.eva.mpg.de/lingua/pdf/Glossing-Rules.pdf}{Leipzig
    Glossing Rules}.
\item Name the language on the first line of the example, with its
  ISO 639-3 code in square brackets: Czech \texttt{[ces]}, English
  \texttt{[eng]}, German \texttt{[deu]}.  Look yours up if you are
  not sure.
\end{itemize}

Cite the source of any example you did not make up yourself.

%%% ------------------------------------------------------------------
\section{Previous work}
%%% ------------------------------------------------------------------

TODO: What have other people said about this?  You need at least three
references to the linguistic literature.  Read them --- do not cite a
paper you have not opened.

There are two citation commands.  Use \verb|\citet| when the author is
part of your sentence, and \verb|\citep| when the citation sits in
brackets at the end:

\begin{quote}
  \citet{shieber1985} argues that Swiss German is not context-free.\\
  Swiss German is not context-free \citep{shieber1985}.
\end{quote}

Give a page number when you point at a specific claim, like this:
\citep[p.~337]{shieber1985}.  A similar argument had already been made
for Dutch \citep{bresnan-etal1982}.

%%% ------------------------------------------------------------------
\section{Analysis}
%%% ------------------------------------------------------------------

TODO: Then do \emph{one} of the following.

\paragraph{Either a formal analysis.}
Work out part of the problem as explicitly as possible, preferably in the framework we have been using \citep{sag-etal2003}.  You do not have to solve all of it --- take a piece small enough that you can be precise about it, and say what still needs to be explained.

\paragraph{some examples} Delete whichever of these you do not need.  There are fuller examples, including type hierarchies and AVMs embedded in trees, in the \href{https://www.overleaf.com/read/hnjksbxgbhtq}{short
  guide to formatting HPSG with LaTeX}.

\begin{center}
\begin{forest}
for tree={
    if n children=0{font=\itshape, tier=terminal}{},
}
  [S [NP [N [Kim]]]
     [VP [V [left]]]]
\end{forest}
\end{center}

\begin{center}
\avm{[ \type*{word} \\
       head & \type{verb} \\
       val  & [ comps & < NP > ] ]}
\end{center}

\paragraph{Or a survey of speakers' grammaticality judgements.}
Ask at least five speakers, report how many accepted each example, and
say what you asked them and how.  Raw counts are fine; you are not
expected to do statistics.  Say clearly where your speakers came from,
because it matters for what the numbers mean.

%%% ------------------------------------------------------------------
\section{Conclusion}
%%% ------------------------------------------------------------------

TODO: What did you find?  Link back to the puzzle you opened with.
End with what you showed, and possibly some ides for future work.

%%% ------------------------------------------------------------------
\subsection*{Contribution statement}
%%% ------------------------------------------------------------------

TODO: Say who did what.  A couple of sentences is enough, for example:
``A.\ collected and glossed the data and wrote Sections 2 and 4;
B.\ did the literature review and wrote Sections 1 and 3; both authors
revised the whole paper.''

%%% ------------------------------------------------------------------

\bibliography{assignment}

\end{document}
